\frontchapter{The Book at a Glance}

Shelley divides two thousand years into nine \emph{ages}, because in each of them the conditions of the
church's life changed. Every age opens with a short paragraph that carries the plot; the chapters then
take up one issue each. Read in a row, the nine introductions give the whole story.

\age{1}{The Age of Jesus and the Apostles}{6 BC--AD 70 \textperiodcentered\ chs.~1--2}
{Jesus of Nazareth starts a renewal movement inside Judaism. After his crucifixion his teaching spreads
around the Mediterranean, and Paul's message of salvation as God's gift for all people turns a Palestinian
Jewish movement into a universal faith.}

\age{2}{The Age of Catholic Christianity}{70--312 \textperiodcentered\ chs.~3--8}
{The faith spreads through the Roman Empire and probably as far as India. Under ridicule and persecution
Christians call their movement \emph{catholic}: universal in reach and true in teaching. It is marked by
a universal vision, orthodox belief and the leadership of bishops.}

\age{3}{The Age of the Christian Roman Empire}{312--590 \textperiodcentered\ chs.~9--16}
{With Constantine the church moves from the catacombs to the palace and ends the fourth century as the
established religion, responsible for a whole society. Doctrine is defined and structures grow. Some
see a victory, the monks see a captivity; when the western empire falls, it is the monks who win the
barbarians.}

\age{4}{The Christian Middle Ages}{590--1517 \textperiodcentered\ chs.~17--23}
{Out of the ruins of the western empire the church builds a new order, Europe, around the idea of
Christendom, the union of empire and church. Papal power climbs to its height under Innocent III and
is then corrupted, while reformers call ever more loudly for change.}

\age{5}{The Age of the Reformation}{1517--1648 \textperiodcentered\ chs.~24--31}
{Reform breaks the papal leadership of the West and produces four Protestant traditions: Lutheran,
Reformed, Anabaptist and Anglican. Rome recovers its fervor under the Jesuits, wars of religion ravage
Europe, and the denominational idea of the church points to a new way of living with division.}

\age{6}{The Age of Reason and Revival}{1648--1789 \textperiodcentered\ chs.~32--35}
{Science and human reason replace the Christian faith as the foundation of Western culture. Many
Protestants answer not with argument but with the experience of conversion; evangelical Christianity
spreads by preaching, and Christians learn to accept religious freedom.}

\age{7}{The Age of Progress}{1789--1914 \textperiodcentered\ chs.~36--41}
{The French Revolution gives the common people new hopes, the cities fill with the poor, and evolution
raises new questions about humanity. Divided over the answers and no longer carried by the state,
Protestants turn to voluntary societies for social work and foreign missions.}

\age{8}{The Age of Ideologies}{1914--1989 \textperiodcentered\ chs.~42--45}
{Communism, Nazism and Americanism struggle for the world, and a new paganism appeals to economics,
race and individual rights. Christians suffer, think and unite in new ways: Protestants in the
ecumenical movement, Roman Catholics in the renewal of their church.}

\age{9}{The Age of Global Expansion and Relocation}{1900-- \textperiodcentered\ chs.~46--48}
{More people have become Christians in the last hundred years than in any earlier period, most of them
south of the Equator, while the old heartlands in Europe and North America decline. Whether the new
centers stay faithful, and whether the West revives, is left to the future.}

\par\addvspace{1.2\baselineskip}
\noindent The book ends with an epilogue in two voices, the source notes, and a list of the popes from
Leo~I to the present.
